\section{Periodic orbital response and the Z-vector}
\label{sec:orbital-response}

\subsection{Response of the orthonormality constraint}

Under a nuclear displacement $R_{Ax}$, the AO overlap changes, so the MO
coefficient derivative is not merely a unitary rotation.  Differentiating
Eq.~\eqref{eq:mo-orthonormality} gives
\begin{equation}
  C^{\vk\dagger}S^{\vk}C^{\vk,Ax}
  +C^{\vk,Ax\dagger}S^{\vk}C^{\vk}
  +C^{\vk\dagger}S^{\vk,Ax}C^{\vk}=0.
  \label{eq:orthonormality-response}
\end{equation}
Defining $U^{\vk,Ax}=C^{\vk\dagger}S^{\vk}C^{\vk,Ax}$ yields
\begin{equation}
  U^{\vk,Ax}+U^{\vk,Ax\dagger}
  =-S_{\mathrm{MO}}^{\vk,Ax},
  \qquad
  S_{\mathrm{MO}}^{\vk,Ax}=C^{\vk\dagger}S^{\vk,Ax}C^{\vk}.
  \label{eq:u-overlap-constraint}
\end{equation}
The occupied--occupied and virtual--virtual gauges may be chosen, whereas the
occupied--virtual block is determined by the HF response equation.

\subsection{The CPHF equation}

Differentiating the HF stationarity condition $F_{ai}^{\vk}=0$ with respect to
a nuclear coordinate gives
\begin{equation}
  \sum_{\vk'bj}A_{ai\vk,bj\vk'}U_{bj}^{\vk',Ax}
  =-B_{ai}^{\vk,Ax}.
  \label{eq:cphf-equation}
\end{equation}
Here $A$ is the periodic restricted HF orbital Hessian and $B^{Ax}$ is the
nuclear perturbation source.  In canonical-orbital notation, the Hessian
action has the schematic form
\begin{equation}
  (AU)_{ai}^{\vk}
  =(\varepsilon_a^{\vk}-\varepsilon_i^{\vk})U_{ai}^{\vk}
  +\left[C_v^{\vk\dagger}\,\delta G_{\mathrm{HF}}[D^{(1)}(U)]^{\vk}
  C_o^{\vk}\right]_{ai}.
  \label{eq:hf-hessian-action}
\end{equation}
Since $\delta G_{\mathrm{HF}}=\delta J-\tfrac12\delta K$, applying the selected
FFTDF/GDF/RSDF $J/K$ provider to a transition density is the core AX
operation.

The nuclear source contains one-electron, two-electron, and overlap terms,
\begin{equation}
  B_{ai}^{\vk,Ax}
  =\left[C_v^{\vk\dagger}
    \left(h^{\vk,Ax}+G_{\mathrm{HF}}^{\vk,Ax}
          -\varepsilon_i^{\vk}S^{\vk,Ax}\right)
    C_o^{\vk}\right]_{ai}
  +B_{ai,\mathrm{basis}}^{\vk,Ax}.
  \label{eq:cphf-source}
\end{equation}
$G_{\mathrm{HF}}^{Ax}$ is the integral derivative at fixed density.  The final
term may be redistributed among other terms depending on the chosen AO-row or
AO-column convention; production code and its documentation must use the same
convention.

\subsection{The adjoint Z-vector}

Solving Eq.~\eqref{eq:cphf-equation} independently for all
$3N_{\mathrm{atom}}$ displacements is expensive.  Collect the target orbital
gradient into a vector $g$ and solve the adjoint equation once,
\begin{equation}
  A^{\Herm}z=-g.
  \label{eq:zvector-equation}
\end{equation}
The chain-rule contribution then becomes
\begin{equation}
  g_{Ax}^{\mathrm{resp}}
  =z^{\Herm}B^{Ax}.
  \label{eq:zvector-contraction}
\end{equation}
The signs in Eqs.~\eqref{eq:zvector-equation} and
\eqref{eq:zvector-contraction} follow from defining the CPHF equation as
$AU=-B$.  If the definition of $B$ or $g$ changes, both signs must be changed
consistently.

\subsection{The real-linear problem at complex k-points}

For a complex orbital rotation $U=X+iY$, the Hermitian first-order density is
\begin{equation}
  D^{(1),\vk}(U)
  =2\left(C_v^{\vk}U^{\vk}C_o^{\vk\dagger}
          +C_o^{\vk}U^{\vk\dagger}C_v^{\vk\dagger}\right).
  \label{eq:first-order-density}
\end{equation}
Because $U^*$ also appears, the map $U\mapsto D^{(1)}$ is real-linear rather
than complex-linear.  The solver coordinate is therefore realified as
\begin{equation}
  u_R=\begin{pmatrix}\operatorname{Re}U\\\operatorname{Im}U\end{pmatrix},
  \label{eq:realified-coordinate}
\end{equation}
and the equations solved are
\begin{equation}
  A_Ru_R=-B_R,
  \qquad A_R^Tz_R=-g_R.
  \label{eq:realified-zvector}
\end{equation}
The fact that Gamma-point matrices can be real is merely a simplifying special
case of this general formulation.

\subsection{k-point normalization and momentum structure}

A nuclear perturbation preserves primitive-cell periodicity and therefore
generates a VO response within each $\vk$.  The $J/K$ Hessian action, however,
couples density responses from all k-points.  Under the equal-weight
convention, both AX and the final scalar contraction must contain exactly one
factor $1/N_k$.  Adding another factor of two without checking whether the
transition density already contains the occupation factor may accidentally
work at Gamma while breaking sampled-k consistency.

\subsection{Gauge freedom, small gaps, and residuals}

Unitary rotations within the occupied or virtual subspaces leave the density
unchanged.  For degenerate orbitals, the subspace-invariant density response,
not an individual $U_{ai}$, is the physical quantity.  The solver operates only
in the VO space and must not silently regularize small denominators.  At a
minimum, the report records
\begin{equation}
  r_{\max}=\|A_R^Tz_R+g_R\|_\infty
  \label{eq:zvector-residual}
\end{equation}
together with the response dimension and a condition estimate.  Even good
gradient--FD agreement is insufficient evidence if $r_{\max}$ is large,
because the result may depend on cancellation for a particular perturbation.
